% !TEX root = ../../ctfp-print.tex

\lettrine[lhang=0.17]{Y}{ou've seen how to model} 类型和纯函数作为范畴的方法。我也提过存在一种在范畴论中建模副作用（非纯函数）的方式。让我们来看其中一个例子：带有日志记录功能的函数。这种在命令式语言中通常通过修改全局状态实现的功能，例如：

\begin{snip}{cpp}
string logger;

bool negate(bool b) {
    logger += "Not so! ";
    return !b;
}
\end{snip}
你知道这不是纯函数，因为其记忆化版本无法产生日志。这个函数具有\newterm{副作用}(side effects)。

现代编程中我们尽量避免使用全局可变状态——这主要是由于并发带来的复杂性。你也绝不会在库函数中编写此类代码。

幸运的是，我们可以使这个函数变得纯净。只需显式传递日志即可。让我们通过添加字符串参数，并将常规输出与更新后的日志配对来实现：

\begin{snip}{cpp}
pair<bool, string> negate(bool b, string logger) {
    return make_pair(!b, logger + "Not so! ");
}
\end{snip}
这个纯函数没有任何副作用，每次用相同参数调用都会返回相同结果，必要时可以记忆化。然而考虑到日志的累积特性，你需要记忆化所有可能导致特定调用的历史记录。例如以下调用都需要独立的记忆化条目：

\begin{snip}{cpp}
negate(true, "It was the best of times. ");
\end{snip}
和

\begin{snip}{cpp}
negate(true, "It was the worst of times. ");
\end{snip}
等等。

这对库函数接口来说也不够友好。调用者虽然可以忽略返回类型的字符串（这不是大问题），但必须强制传入字符串参数，这会带来不便。

有没有更优雅的实现方式？能否分离关注点？在这个简单示例中，\code{negate}的主要功能是转换布尔值，日志记录是次要功能。诚然，记录的消息与函数相关，但将消息聚合为完整日志是独立需求。我们仍希望函数产生字符串，但希望解除其生成日志的责任。折中方案如下：

\begin{snip}{cpp}
pair<bool, string> negate(bool b) {
    return make_pair(!b, "Not so! ");
}
\end{snip}
关键思想是日志将在\emph{函数调用之间}聚合。

为了展示具体实现，我们换一个更现实的例子。假设有两个函数：一个将字符串转为大写：

\begin{snip}{cpp}
string toUpper(string s) {
    string result;
    int (*toupperp)(int) = &toupper; // toupper 被重载
    transform(begin(s), end(s), back_inserter(result), toupperp);
    return result;
}
\end{snip}
另一个将字符串按空白字符分割为字符串向量：

\begin{snip}{cpp}
vector<string> toWords(string s) {
    return words(s);
}
\end{snip}
实际工作由辅助函数\code{words}完成：

\begin{snip}{cpp}
vector<string> words(string s) {
    vector<string> result{""};
    for (auto i = begin(s); i != end(s); ++i)
    {
        if (isspace(*i))
            result.push_back(""); 
        else
            result.back() += *i;
    }
    return result;
}
\end{snip}
我们需要修改\code{toUpper}和\code{toWords}，使其在常规返回值基础上附带消息字符串。

\begin{figure}[H]
  \centering
  \includegraphics[width=0.3\textwidth]{images/piggyback.jpg}
\end{figure}
\noindent
我们将"修饰"这些函数的返回值。通过定义模板\code{Writer}实现通用封装，它包装一个二元组，第一个分量为任意类型\code{A}的值，第二个分量为字符串：

\begin{snip}{cpp}
template<class A>
using Writer = pair<A, string>;
\end{snip}
修饰后的函数如下：

\begin{snip}{cpp}
Writer<string> toUpper(string s) {
    string result;
    int (*toupperp)(int) = &toupper;
    transform(begin(s), end(s), back_inserter(result), toupperp);
    return make_pair(result, "toUpper "); 
}

Writer<vector<string>> toWords(string s) { 
    return make_pair(words(s), "toWords ");
}
\end{snip}
我们需要组合这两个函数形成新的修饰函数，实现字符串大写转换和分词操作，同时生成操作日志。具体实现如下：

\begin{snip}{cpp}
Writer<vector<string>> process(string s) {
    auto p1 = toUpper(s);
    auto p2 = toWords(p1.first);
    return make_pair(p2.first, p1.second + p2.second);
}
\end{snip}
我们达成了目标：日志聚合不再是各个函数的职责。它们各自生成消息，外部再将这些消息拼接成完整日志。

想象整个程序都采用这种风格编写，这将是重复且易出错的噩梦。但作为程序员，我们知道如何应对重复代码：抽象！但这不是普通的抽象——我们需要抽象\newterm{函数组合}(function composition)本身。而组合正是范畴论的本质，在编写更多代码前，让我们从范畴论角度分析这个问题。

\section{Writer 范畴}

通过修饰一组函数的返回类型来附加额外功能的想法非常有价值。我们会看到更多这类示例。起点是我们熟悉的类型与函数构成的范畴。保持类型作为对象不变，但将态射重新定义为修饰后的函数。

例如，考虑修饰从\code{int}到\code{bool}的函数\code{isEven}。我们将其转化为由修饰函数表示的态射。关键在于这个态射仍然被视为对象\code{int}到\code{bool}的箭头，尽管修饰函数返回的是二元组：

\begin{snip}{cpp}
pair<bool, string> isEven(int n) {
    return make_pair(n % 2 == 0, "isEven ");
}
\end{snip}
根据范畴律，我们应该能将这个态射与从对象\code{bool}到其他对象的态射组合。特别是可以与之前的\code{negate}组合：

\begin{snip}{cpp}
pair<bool, string> negate(bool b) {
    return make_pair(!b, "Not so! ");
}
\end{snip}
显然，我们不能像普通函数那样组合这两个态射，因为输入/输出不匹配。它们的组合应该如下：

\begin{snip}{cpp}
pair<bool, string> isOdd(int n) {
    pair<bool, string> p1 = isEven(n);
    pair<bool, string> p2 = negate(p1.first);
    return make_pair(p2.first, p1.second + p2.second);
}
\end{snip}
这是我们构建的新范畴中态射组合的具体规则：

\begin{enumerate}
  \tightlist
  \item
        执行对应第一个态射的修饰函数
  \item
        提取结果二元组的第一个分量，传递给对应第二个态射的修饰函数
  \item
        拼接第一个结果的第二个分量（字符串）与第二个结果的第二个分量（字符串）
  \item
        返回新二元组，组合最终结果的第一个分量与拼接后的字符串
\end{enumerate}

若要在C++中将这种组合抽象为高阶函数，我们需要使用模板，参数化三个对应范畴中三个对象的类型。它接受两个按规则可组合的修饰函数，返回第三个修饰函数：

\begin{snip}{cpp}
template<class A, class B, class C>
function<Writer<C>(A)> compose(function<Writer<B>(A)> m1,
                               function<Writer<C>(B)> m2)
{
    return [m1, m2](A x) {
        auto p1 = m1(x);
        auto p2 = m2(p1.first);
        return make_pair(p2.first, p1.second + p2.second); 
    };
}
\end{snip}
现在我们可以回到前面的例子，使用这个新模板组合\code{toUpper}和\code{toWords}：

\begin{snip}{cpp}
Writer<vector<string>> process(string s) { 
    return compose<string, string, vector<string>>(toUpper, toWords)(s);
}
\end{snip}
传递类型参数给\code{compose}模板仍有较多冗余。只要使用支持泛型lambda和返回类型推导的C++14编译器，就可以简化（此代码归功于Eric Niebler）：

\begin{snip}{cpp}
auto const compose = [](auto m1, auto m2) { 
    return [m1, m2](auto x) { 
        auto p1 = m1(x);
        auto p2 = m2(p1.first);
        return make_pair(p2.first, p1.second + p2.second);
    };
};
\end{snip}
采用新定义后，\code{process}的实现简化为：

\begin{snip}{cpp}
Writer<vector<string>> process(string s) {
    return compose(toUpper, toWords)(s);
}
\end{snip}
但我们尚未完成。虽然定义了组合规则，但单位元态射是什么？它们不是普通的恒等函数！它们是从类型A到自身的态射，即形式为：

\begin{snip}{cpp}
Writer<A> identity(A);
\end{snip}
它们必须满足组合的单位元性质。观察我们的组合定义可知，恒等态射应该无改动传递参数，并向日志贡献空字符串：

\begin{snip}{cpp}
template<class A> Writer<A> identity(A x) {
    return make_pair(x, "");
}
\end{snip}
你可以验证这个定义确实满足范畴律。组合的结合性成立是因为：第一个分量的运算是普通函数组合（满足结合性），第二个分量的拼接也满足结合性。

敏锐的读者可能注意到，这个构造可以轻松推广到任意幺半群，不限于字符串。我们只需在\code{compose}中使用\code{mappend}，在\code{identity}中使用\code{mempty}（代替\code{+}和\code{""}）。实际上没有理由限制日志类型仅为字符串。优秀的库设计者应该识别使库工作的最小约束条件——此处日志库唯一需要的要求是日志具有幺半群性质。

\section{Haskell中的Writer}

在Haskell中实现相同的功能会更加简洁，而且我们还能获得编译器的更多支持。让我们先定义\code{Writer}类型：

\src{snippet01}
这里我只是定义了一个类型别名，相当于C++中的\code{typedef}（或\code{using}）。类型\code{Writer}由一个类型变量\code{a}参数化，等价于\code{a}和\code{String}的有序对。对于有序对的语法表示非常简洁：只需用括号包裹两个元素，中间用逗号分隔即可。

我们的态射是从任意类型到某个\code{Writer}类型的函数：

\src{snippet02}
我们将把组合操作定义为一个特殊的中缀运算符，有时被称为"fish"：

\src{snippet03}
这是一个接受两个函数作为参数并返回新函数的运算符。第一个参数类型为\code{(a -> Writer b)}，第二个参数类型为\code{(b -> Writer c)}，结果类型为\code{(a -> Writer c)}。

这个中缀运算符的具体定义如下——其中\code{m1}和\code{m2}分别出现在fish符号的左右两侧：

\src{snippet04}
结果是一个以\code{x}为参数的λ函数。在Haskell中λ函数用反斜杠表示——可以理解为希腊字母$\lambda$被截断了一条腿。

\code{let}表达式允许我们声明辅助变量。在这里，调用\code{m1}的结果通过模式匹配分解为变量对\code{(y, s1)}；而用第一个模式得到的\code{y}作为参数调用\code{m2}的结果，又被匹配为\code{(z, s2)}。

在Haskell中模式匹配有序对是很常见的做法，这与我们在C++中使用访问器的方式不同。除此之外，这两种实现之间存在相当直接的对应关系。

\code{let}表达式的整体返回值由其\code{in}子句指定：这里返回的是一个有序对，第一个分量是\code{z}，第二个分量是两个字符串\code{s1++s2}的连接结果。

我们还需要为这个范畴定义单位态射，不过出于后文才会阐明的原因，我将其命名为\code{return}：

\src{snippet05}
为了完整性，让我们看看装饰函数\code{upCase}和\code{toWords}的Haskell版本：

\src{snippet06}
这里的\code{map}函数对应C++中的\code{transform}。它将字符函数\code{toUpper}应用到字符串\code{s}上。辅助函数\code{words}已经在标准预置库Prelude中定义。

最后，这两个函数的组合可以通过fish运算符来完成：

\src{snippet07}

\section{克莱斯利范畴}

你或许已经猜到，这个范畴并不是我临时编造的。它属于所谓的克莱斯利范畴（Kleisli category）--- 一种基于单子（monad）构建的范畴。虽然我们尚未准备好深入讨论单子理论，但我希望通过这个例子让你初步了解单子能实现的功能。就当前有限的目标而言，克莱斯利范畴的对象是底层编程语言中的类型。从类型 $A$ 到类型 $B$ 的态射（morphism）是形如 $A \to B'$ 的函数，其中 $B'$ 是通过特定"装饰"（embellishment）操作从类型 $B$ 派生而来。每个克莱斯利范畴都定义了特有的态射组合规则，以及与该组合操作相容的恒等态射（identity morphism）。（后文将揭示，此处模糊使用的"装饰"一词实际上对应于范畴论中的内函子（endofunctor）概念。）

本章实例所采用的具体单子被称为\newterm{writer monad}（记录单子），主要用于函数执行过程的日志记录或跟踪。它也是将副作用嵌入纯计算的一种更通用机制的典型示例。此前我们曾用集合范畴对编程语言的类型和函数进行建模（照常忽略底值问题），而今我们已将此模型拓展至一个略有不同的范畴：在这个新范畴中，态射由经过装饰的函数表示，其组合操作不再局限于简单地将一个函数的输出传递给另一个函数作为输入。我们获得了额外的操作自由度——组合本身的构造方式。事实证明，正是这种自由度使得我们可以为命令式语言中传统上依赖副作用实现的程序赋予简洁的指称语义学解释。

\section{挑战}

一个并非为其参数的所有可能取值都定义的函数被称为部分函数（partial function）。从数学意义上说它并不是真正的函数，因此无法适配标准的范畴论模型。然而，我们可以通过返回修饰类型\code{optional}的函数来表示它：

\begin{snip}{cpp}
template<class A> class optional {
    bool _isValid;
    A _value;
public: 
    optional()    : _isValid(false) {}
    optional(A v) : _isValid(true), _value(v) {}
    bool isValid() const { return _isValid; }
    A value() const { return _value; }
};
\end{snip}
例如，以下是修饰函数\code{safe\_root}的实现：

\begin{snip}{cpp}
optional<double> safe_root(double x) {
    if (x >= 0) return optional<double>{sqrt(x)}; 
    else return optional<double>{};
}
\end{snip}
现在提出以下挑战：

\begin{enumerate}
  \tightlist
  \item
        构造部分函数的Kleisli范畴（定义复合运算与单位元）。
  \item
        实现修饰函数\code{safe\_reciprocal}，当其参数非零时返回有效的倒数。
  \item
        组合函数\code{safe\_root}与\code{safe\_reciprocal}以实现\code{safe\_root\_reciprocal}，使其能够尽可能计算表达式\code{sqrt(1/x)}。
\end{enumerate}